\documentclass{article}
\usepackage[utf8]{inputenc}
\usepackage[T1]{fontenc}
\usepackage{amsmath,amssymb}
\usepackage{graphicx}
\usepackage{booktabs}
\usepackage{hyperref}
\usepackage{natbib}
\usepackage[margin=1in]{geometry}

\title{Benchmarking DNA Foundation Models for Genomic Sequence Classification}
\author{[Author Names]}
\date{}

\begin{document}
\maketitle

% ============================================================
\begin{abstract}
The rapid proliferation of DNA foundation language models---each evaluated on
different tasks with incompatible experimental setups---has left practitioners
without evidence-based guidance for model selection. We present a unified
benchmark evaluating five DNA foundation models (DNABERT-2, Nucleotide
Transformer v2, HyenaDNA, Caduceus-Ph, and GROVER) as zero-shot feature
extractors across 57 classification datasets spanning four task categories:
genomic sequence classification, gene expression prediction, variant effect
prediction, and topologically associating domain (TAD) boundary recognition.
Using random forest classifiers trained on frozen embeddings with five-fold
cross-validation, we find that no single model dominates. Caduceus-Ph achieves
the highest AUC on human promoter detection (0.987) and coding region
identification (0.974), while DNABERT-2 leads on splice site prediction (0.906).
Mean token pooling consistently outperforms the default summary token strategy,
improving AUC by 4--9\% across models. For variant effect prediction, the
Nucleotide Transformer v2 achieves the best performance on pathogenic SNP
classification (AUC 0.732), yet general-purpose foundation models substantially
underperform specialized architectures on quantitative trait locus (QTL)
prediction, where AlphaGenome achieves an AUC of 0.803 on eQTL versus 0.649 for
the best foundation model. A controlled ablation demonstrates that multi-species
pre-training data improves cross-species transfer on 14 of 49 datasets. These
findings establish practical guidelines for task-specific model selection and
reveal a fundamental specialization gap in current genomic language models.
\end{abstract}

% ============================================================
\section{Introduction}
\label{sec:intro}

The sequencing of the human genome~\citep{lander_2001_human} and subsequent
large-scale annotation efforts~\citep{encode_2012, roadmap_2015} have generated
vast catalogs of genomic elements, yet deciphering the functional code within
DNA sequences remains a central challenge. The application of large-scale
pre-trained language models to biological sequences has emerged as a
transformative paradigm~\citep{ji_2021_dnabert, elnaggar_2022_prottrans,
rives_2021_esm}. Inspired by the success of masked language modeling in natural
language processing~\citep{devlin_2019_bert, vaswani_2017_attention}, recent work
has produced a family of DNA foundation models that learn contextualized
representations from large genomic corpora~\citep{zhou_2024_dnabert2,
dallatorre_2024_nucleotide, nguyen_2023_hyenadna, schiff_2024_caduceus,
sanabria_2024_grover}.

These models span diverse architectural paradigms---from BERT-style transformers
with byte-pair encoding~\citep{zhou_2024_dnabert2, sennrich_2016_bpe} and
multi-species pre-training~\citep{dallatorre_2024_nucleotide}, to sub-quadratic
alternatives based on Hyena operators~\citep{nguyen_2023_hyenadna,
poli_2023_hyena} and selective state space models~\citep{schiff_2024_caduceus,
gu_2024_mamba}. Each model has been evaluated on a different subset of genomic
tasks using incompatible benchmarks, making fair comparison across models
impossible. Existing benchmark suites such as the Genome Understanding Evaluation
(GUE)~\citep{zhou_2024_dnabert2} and
GenomicBenchmarks~\citep{gresova_2023_genomicbenchmarks} provide standardized
datasets but do not systematically compare the full set of contemporary models
under identical conditions.

Despite growing interest in applying machine learning to genomic
sequences~\citep{ji_2021_survey_dna_ml}, several critical questions remain
unanswered. First, which model performs best
for a given genomic task type? Second, how sensitive are model comparisons to the
choice of embedding pooling strategy---a hyperparameter that has received little
attention in the genomic modeling literature? Third, can general-purpose
foundation models match or replace specialized architectures such as
Enformer~\citep{avsec_2021_enformer} and Sei~\citep{chen_2022_sei} on tasks
requiring variant-level or expression-level prediction?

In this work, we address these gaps through a comprehensive, controlled benchmark
study. Our contributions are as follows:
\begin{itemize}
\item We conduct a unified evaluation of five DNA foundation models across 57
      datasets spanning four task categories, using identical evaluation protocols
      and downstream classifiers.
\item We demonstrate that mean token pooling consistently outperforms the default
      summary token embedding by 4--9\% AUC, with statistical significance
      confirmed by DeLong's test~\citep{delong_1988_comparing} on 35--42 of 52
      binary datasets per model.
\item We provide task-specific model rankings showing that no single model
      dominates: model choice should be guided by the specific genomic task.
\item We identify a fundamental specialization gap: general foundation models
      substantially underperform specialized architectures on QTL prediction.
\item We present a controlled ablation showing that multi-species pre-training
      data improves cross-species genomic transfer on 14 of 49 datasets.
\end{itemize}

% ============================================================
\section{Related Work}
\label{sec:related}

\subsection{DNA Foundation Models}

The original DNABERT~\citep{ji_2021_dnabert} applied BERT-style masked language
modeling to k-mer tokenized DNA sequences, demonstrating that pre-trained
representations transfer effectively to promoter, splice site, and transcription
factor binding site prediction. DNABERT-2~\citep{zhou_2024_dnabert2} replaced
k-mer tokenization with byte-pair encoding~\citep{sennrich_2016_bpe} and
introduced ALiBi positional encoding~\citep{press_2022_alibi}, achieving
comparable performance with 21$\times$ fewer parameters. The Nucleotide
Transformer~\citep{dallatorre_2024_nucleotide} scaled pre-training to 850
species genomes with up to 2.5 billion parameters.

Beyond transformer architectures, HyenaDNA~\citep{nguyen_2023_hyenadna} employed
Hyena operators~\citep{poli_2023_hyena} to achieve sub-quadratic scaling with
context lengths up to one million nucleotides. Caduceus~\citep{schiff_2024_caduceus}
introduced bi-directional and reverse-complement equivariant extensions of the
Mamba state space model~\citep{gu_2024_mamba}. GROVER~\citep{sanabria_2024_grover}
applied a custom next-k-mer prediction objective to learn genomic grammar.

\subsection{Genomic Sequence Classification Benchmarks}

The GUE benchmark~\citep{zhou_2024_dnabert2} comprises 36 datasets across nine
task categories, while GenomicBenchmarks~\citep{gresova_2023_genomicbenchmarks}
provides eight curated datasets for promoter, enhancer, and open chromatin
classification. However, these benchmarks were designed to evaluate individual
models and do not enforce consistent evaluation protocols across all contemporary
architectures.

\subsection{Specialized Genomic Prediction Models}

For gene expression prediction, Enformer~\citep{avsec_2021_enformer} integrates
information from up to 200 kilobases of flanking sequence, substantially
extending the receptive field of earlier convolutional
models~\citep{kelley_2018_basenji, kelley_2020_basenji2}. For variant effect
prediction, Sei~\citep{chen_2022_sei} provides a sequence-based map of 21,907
regulatory activities, while tools such as CADD~\citep{kircher_2014_cadd},
AlphaMissense~\citep{cheng_2023_accurate}, and deep generative
models~\citep{frazer_2021_disease} combine multiple evidence types for
pathogenicity scoring. Recent work has demonstrated that DNA language models
can serve as effective variant effect predictors~\citep{benegas_2023_dnabert_survey},
though a systematic comparison against specialized architectures across multiple
variant types has been lacking.

\subsection{Representation Learning and Pooling Strategies}

In NLP, the choice between CLS token and mean token pooling for sentence
embeddings has been extensively studied~\citep{devlin_2019_bert}, with mean
pooling often outperforming CLS on downstream tasks. In genomics, this design
choice has received little systematic attention, with most models defaulting
to their pre-trained summary token strategy.

% ============================================================
\section{Methods}
\label{sec:methods}

\subsection{Models Under Evaluation}

We evaluate five DNA foundation models spanning four architectural families
(Table~\ref{tab:models}). As baselines, we include three specialized models
(Enformer, Sei, and AlphaGenome) and a task-specific convolutional neural
network (CNN) with three layers of 64, 128, and 256 channels trained
end-to-end on each dataset.

\begin{table}[t]
\centering
\caption{DNA foundation models evaluated in this benchmark. All models are used
in zero-shot mode with frozen weights.}
\label{tab:models}
\begin{tabular}{lccccc}
\toprule
Model & Params & Architecture & Emb.\ Dim & Max Input & Pre-training \\
\midrule
DNABERT-2 & 117M & BERT + ALiBi & 768 & No limit & 135 species \\
NT-v2 & 500M & BERT & 1024 & 12K nt & 850 species \\
HyenaDNA & 30M & Hyena & 256 & 1M nt & Human only \\
Caduceus-Ph & $\sim$35M & BiMamba (SSM) & 256 & 131K nt & Multi-species \\
GROVER & 117M & Transformer & 768 & $\sim$2K nt & Multi-species \\
\bottomrule
\end{tabular}
\end{table}

\subsection{Task Taxonomy and Datasets}

We organize 57 datasets into four task categories:
(1)~\textbf{Genomic sequence classification}, including human genome tasks
(splice site, promoter, enhancer, coding region, transcription factor binding
site detection), multi-species tasks (cross-organism classification, plant
promoter detection), and epigenetic modification prediction (5mC methylation,
histone marks in yeast and human);
(2)~\textbf{Gene expression prediction} from promoter-proximal sequences at
short (6K~bp) and long (196K~bp) context lengths;
(3)~\textbf{Variant effect prediction}, including pathogenic versus common SNP
classification using ClinVar~\citep{landrum_2018_clinvar} and
gnomAD~\citep{karczewski_2020_gnomad} annotations, with common variants defined
from the 1000 Genomes Project~\citep{auton_2015_1000genomes}, and quantitative
trait locus (eQTL and ipaQTL) prediction; and
(4)~\textbf{TAD boundary recognition}.

\subsection{Evaluation Protocol}

All foundation models are used in zero-shot mode: model weights are frozen, and
embeddings are extracted for each input DNA sequence. We train a random forest
classifier~\citep{breiman_2001_random} on these embeddings using five-fold
cross-validation with hyperparameter optimization. For binary classification
tasks, we report the area under the receiver operating characteristic curve
(AUC). For multi-class tasks, we report accuracy. For regression tasks (gene
expression), we report Pearson correlation and mean squared error (MSE).

\subsection{Pooling Strategy Comparison}

For each model, we compare three embedding pooling strategies: (1)~the default
summary/CLS token, (2)~mean pooling across all output token embeddings, and
(3)~max pooling. Statistical significance of AUC differences between pooling
strategies is assessed using DeLong's test~\citep{delong_1988_comparing} with a
significance threshold of $p < 0.01$.

\subsection{Pre-training Data Ablation}

To assess the impact of pre-training data diversity, we retrained HyenaDNA---the
only model originally trained solely on the human genome---on a multi-species
corpus. We compared the original and multi-species variants across all 49
applicable binary classification datasets using paired Wilcoxon signed-rank
tests.

% ============================================================
\section{Results}
\label{sec:results}

\subsection{Sequence Classification}
\label{sec:results_classification}

Table~\ref{tab:classification} presents AUC results across selected genomic
sequence classification tasks using mean token pooling (the recommended strategy;
see Section~\ref{sec:results_pooling}).

\begin{table}[t]
\centering
\caption{Sequence classification AUC (mean token pooling) on selected tasks.
Best result per row in \textbf{bold}. All models used in zero-shot mode with
random forest downstream classifiers.}
\label{tab:classification}
\begin{tabular}{lccccc}
\toprule
Task & DNABERT-2 & NT-v2 & HyenaDNA & Caduceus & GROVER \\
\midrule
\multicolumn{6}{l}{\textit{Human genome}} \\
Splice donor     & \textbf{0.906} & 0.820 & 0.813 & 0.854 & 0.819 \\
Splice acceptor  & \textbf{0.897} & 0.793 & 0.795 & 0.845 & 0.804 \\
Coding region    & 0.944 & 0.929 & 0.941 & \textbf{0.974} & 0.959 \\
Promoter (GM12878) & 0.986 & 0.984 & 0.976 & \textbf{0.987} & 0.984 \\
Enhancer         & \textbf{0.872} & 0.867 & 0.834 & 0.838 & 0.855 \\
TFBS             & 0.838 & 0.832 & 0.830 & \textbf{0.880} & 0.862 \\
\midrule
\multicolumn{6}{l}{\textit{Multi-species}} \\
Arabidopsis TATA & 0.951 & 0.950 & \textbf{0.961} & 0.937 & 0.949 \\
Human vs.\ worm  & 0.980 & 0.979 & 0.950 & \textbf{0.992} & 0.984 \\
\midrule
\multicolumn{6}{l}{\textit{Epigenetic modifications}} \\
Human 5mC        & 0.685 & 0.738 & 0.684 & \textbf{0.783} & 0.744 \\
Yeast H3         & 0.914 & 0.895 & 0.900 & \textbf{0.929} & 0.906 \\
Yeast H4         & \textbf{0.931} & 0.910 & 0.898 & 0.930 & 0.908 \\
\bottomrule
\end{tabular}
\end{table}

No single model dominates across all tasks. Caduceus-Ph achieves the highest AUC
on five of eleven tasks shown, particularly excelling at coding region detection
(0.974), promoter classification (0.987), and epigenetic mark prediction (0.783
on human 5mC). DNABERT-2 leads on splice site prediction (0.906 donor, 0.897
acceptor) and enhancer detection (0.872). HyenaDNA, despite having the fewest
parameters (30M), achieves the best performance on Arabidopsis TATA promoter
classification (0.961), suggesting that its human-only pre-training transfers
effectively to plant genomes when long-range context is available.

\subsection{Impact of Pooling Strategy}
\label{sec:results_pooling}

Switching from the default summary token to mean token pooling yields substantial
and consistent improvements across all models (Figure~\ref{fig:pooling}).
HyenaDNA benefits most (+8.7\% median AUC improvement, IQR 4.6--12.9\%),
followed by NT-v2 (+6.8\%, IQR 3.7--9.6\%) and Caduceus-Ph (+5.9\%, IQR
2.8--9.1\%). DNABERT-2 shows a moderate improvement (+4.0\%, IQR 2.0--5.5\%),
while GROVER is least affected (+1.4\%, IQR 0.7--1.9\%). DeLong's
test~\citep{delong_1988_comparing} confirms significance ($p < 0.01$) on 35--42
of 52 binary datasets per model, indicating that the improvement is robust across
tasks rather than driven by outliers.

\begin{figure}[t]
\centering
\includegraphics[width=0.8\columnwidth]{fig2_pooling.pdf}
\caption{Impact of pooling strategy on model performance. Bars show median AUC
improvement when switching from summary/CLS token to mean token pooling. Error
bars indicate interquartile range across 52 binary classification datasets.}
\label{fig:pooling}
\end{figure}

\subsection{Gene Expression Prediction}
\label{sec:results_expression}

At short context lengths (6K~bp), all foundation models achieve near-identical
performance (Pearson correlation $\approx 0.12$, MSE $\approx 0.235$),
suggesting that promoter-proximal sequence alone provides insufficient signal for
expression prediction regardless of model architecture. At long context (196K~bp),
HyenaDNA (using its 450K-token variant) achieves a Pearson correlation of 0.137
with MSE 0.226, modestly but significantly ($p < 0.001$) outperforming
Enformer~\citep{avsec_2021_enformer} (correlation 0.129, MSE 0.227). This result
demonstrates that HyenaDNA's sub-quadratic architecture enables effective use of
long-range genomic context for expression prediction.

\subsection{Variant Effect Prediction}
\label{sec:results_variant}

Figure~\ref{fig:variant} presents variant effect prediction results across three
task types. For pathogenic versus common SNP classification, the Nucleotide
Transformer v2 achieves the highest AUC (0.732, Cohen's $d = 0.881$), followed
by Caduceus (0.696) and Enformer (0.688). DNABERT-2 performs unexpectedly poorly
(AUC 0.538, Cohen's $d = 0.134$~\citep{cohen_1988_statistical}), suggesting
that its BPE tokenization may not effectively capture single-nucleotide variant
effects.

\begin{figure}[t]
\centering
\includegraphics[width=0.85\columnwidth]{fig3_variant.pdf}
\caption{Variant effect prediction AUC across three task types. General-purpose
foundation models (warm colors) are compared against specialized architectures
(gray: Enformer; red: AlphaGenome). AlphaGenome is not applicable to the
pathogenic SNP task (shown with reduced opacity).}
\label{fig:variant}
\end{figure}

On QTL prediction tasks, specialized models substantially outperform all
foundation models. AlphaGenome achieves AUC 0.803 on eQTL and 0.864 on ipaQTL,
compared to the best foundation model (Caduceus, 0.649 on eQTL; NT-v2, 0.602 on
ipaQTL). Enformer similarly outperforms foundation models on both tasks (0.774
eQTL, 0.692 ipaQTL). This specialization gap indicates that zero-shot
embeddings from general-purpose DNA models do not capture the fine-grained
regulatory signals required for quantitative trait prediction.

\subsection{Pre-training Data Diversity}
\label{sec:results_pretraining}

Retraining HyenaDNA on a multi-species corpus (compared to its original
human-only pre-training) produced significant improvements on 14 of 49 datasets,
with the largest gain on human 5mC methylation prediction (AUC 0.707 $\rightarrow$
0.749, $+5.9\%$). However, three datasets showed decreased performance (human
enhancer, open chromatin, and yeast epigenetic marks), suggesting that
multi-species pre-training introduces trade-offs: while cross-species transfer
improves, some human-specific regulatory signals may be diluted.

\subsection{Computational Efficiency}
\label{sec:results_runtime}

Runtime analysis reveals distinct scaling behaviors. GROVER achieves the lowest
latency on short sequences ($<$2K~nt), while HyenaDNA scales most favorably
to long sequences ($>$100K~nt) due to its sub-quadratic Hyena
operator~\citep{dao_2022_flashattention}. NT-v2
exhibits near-constant latency ($\sim$2.5s) regardless of input length, likely
due to its fixed 12K~nt input window and associated padding overhead.

% ============================================================
\section{Discussion}
\label{sec:discussion}

\subsection{Practical Model Selection Guidelines}

Our results support a task-aware approach to model selection rather than a
one-size-fits-all recommendation. For splice site and enhancer detection tasks,
DNABERT-2 provides the best performance. For broader regulatory element
classification (promoters, coding regions, transcription factor binding sites),
Caduceus-Ph is consistently strongest. For pathogenic variant scoring, NT-v2 is
preferred. For tasks requiring long-range context (gene expression, long-range
regulatory elements), HyenaDNA offers the best scalability-performance trade-off.
In all cases, mean token pooling should be used over default summary token
strategies.

\subsection{The Specialization Gap}

The substantial underperformance of foundation models on QTL prediction relative
to specialized architectures (AlphaGenome, Enformer) reveals a fundamental
limitation of zero-shot transfer from general-purpose DNA language models.
Specialized models for variant effect prediction are trained with explicit
variant-aware objectives that capture allele-specific regulatory perturbations,
while foundation models learn broader sequence-level representations. This gap
suggests that supervised fine-tuning with variant-level labels may be necessary
for foundation models to compete on quantitative trait prediction tasks.

\subsection{Pooling as an Overlooked Hyperparameter}

The consistent 4--9\% AUC improvement from mean token pooling over default
summary tokens highlights an underappreciated design choice. Default CLS or
summary tokens compress an entire sequence into a single representation, which
may discard distributional information that is captured by averaging across all
positions. This finding parallels observations in NLP~\citep{devlin_2019_bert}
and suggests that practitioners should always evaluate mean pooling before
defaulting to model-provided summary tokens.

\subsection{Limitations and Future Work}

This benchmark evaluates models in zero-shot mode with random forest downstream
classifiers~\citep{breiman_2001_random}. Performance rankings may differ under
fine-tuning or with neural downstream heads~\citep{hastie_2009_elements}. Our dataset coverage, while
extensive (57 datasets), is biased toward the human genome and model organisms;
evaluation on clinical variant interpretation~\citep{landrum_2018_clinvar} and
non-model organisms warrants further study. Future work should compare
zero-shot transfer against fine-tuned performance and assess the impact of
model scale within each architecture family.

% ============================================================
\section{Conclusion}
\label{sec:conclusion}

We presented a unified benchmark evaluating five DNA foundation models across 57
genomic classification datasets. Our analysis reveals that no single model
dominates: task-specific model selection is essential for optimal performance.
Mean token pooling provides a universal improvement of 4--9\% AUC over default
strategies. Critically, general-purpose foundation models fail to match
specialized architectures on quantitative trait prediction, with AlphaGenome
(AUC 0.803 on eQTL) substantially outperforming the best foundation model
(0.649). These findings provide actionable guidance for practitioners and
identify the specialization gap as a key challenge for the next generation of
genomic language models.

% ============================================================
\bibliographystyle{plainnat}
\bibliography{references}

\end{document}
